\section{Dasar Matematika: Galois Field \texorpdfstring{$\text{GF}(2^8)$}{GF(2\textasciicircum8)}}
\label{sec:gf-aes}

Operasi \textit{MixColumns} dan konstruksi S-box memakai aritmetika di dalam
$\text{GF}(2^8)$. Sub-bab ini menjelaskan aritmetika tersebut.

\subsection{Medan Berhingga dan Representasi Polinomial}

$\text{GF}(2^m)$ disebut juga medan berhingga biner. $\text{GF}(2^m)$ adalah ruang
vektor berdimensi $m$ pada $\text{GF}(2)$. Setiap elemennya adalah bilangan bulat
dalam representasi biner sepanjang maksimal $m$ bit. Contohnya $\text{GF}(2^4)$
dan $\text{GF}(2^8)$.

String biner $b_{m-1}b_{m-2} \dots b_1b_0$ dengan $b_i \in \{0,1\}$ dapat
dinyatakan sebagai polinom

\[
b_{m-1}x^{m-1} + b_{m-2}x^{m-2} + \dots + b_1x + b_0 .
\]
Jadi setiap $a \in \text{GF}(2^m)$ dapat dinyatakan dalam bentuk polinom tersebut.
Sebagai contoh, $1101$ di dalam $\text{GF}(2^4)$ adalah polinom $x^3 + x^2 + 1$.
Bilangan $01010111$ di dalam $\text{GF}(2^8)$ adalah polinom
$x^6 + x^4 + x^2 + x + 1$.

\subsection{Penjumlahan di \texorpdfstring{$\text{GF}(2^m)$}{GF(2\textasciicircum m)}}

Misalkan $a = (a_{m-1} \dots a_1 a_0)$ dan $b = (b_{m-1} \dots b_1 b_0)$ adalah
dua elemen $\text{GF}(2^m)$. Penjumlahannya dilakukan komponen demi komponen,
\[
a + b = c = (c_{m-1} \dots c_1 c_0), \qquad
c_i = (a_i + b_i) \bmod 2,
\]
dan hasilnya tetap berada di $\text{GF}(2^m)$. Karena koefisien hanya bernilai
0 atau 1, penjumlahan itu sama dengan operasi XOR.

Sebagai contoh, misalkan $a = \code{0D} = x^3 + x^2 + 1$ dan
$b = \code{06} = x^2 + x$. Keduanya berada di $\text{GF}(2^8)$.
\[
a + b = (x^3 + x^2 + 1) + (x^2 + x)
= x^3 + 2x^2 + x + 1 \equiv x^3 + x + 1 \pmod 2.
\]
Hasilnya $x^3 + x + 1 = \code{00001011} = \code{0B}$. Dalam representasi biner,
\begin{center}
\begin{tabular}{@{}r@{}}
\code{00001101} \\
\code{00000110} $\oplus$ \\
\hline
\code{00001011}
\end{tabular}
\end{center}
Jadi $a + b = \code{0B} = a \mathbin{\text{XOR}} b$.

Contoh kedua adalah $a = \code{57} = x^6 + x^4 + x^2 + x + 1$ dan
$b = \code{83} = x^7 + x + 1$. Penjumlahannya menghasilkan
\[
a + b = x^7 + x^6 + x^4 + x^2 + 2x + 2
\equiv x^7 + x^6 + x^4 + x^2 = \code{11010100} = \code{D4}.
\]
Hasil itu sama dengan $\code{57} \oplus \code{83} = \code{D4}$.

\subsection{Perkalian dan Polinom Pereduksi}

Perkalian dua elemen $\text{GF}(2^m)$ dilakukan dengan mengalikan kedua polinom
seperti biasa. Derajat hasilnya dapat mencapai $2m-2$. Hasil itu kemudian
direduksi modulo sebuah \textit{irreducible polynomial} berderajat $m$, sehingga
derajat sisanya kembali di bawah $m$. Dengan demikian hasilnya menjadi elemen
$\text{GF}(2^m)$ lagi.

Algoritma Rijndael beroperasi di dalam $\text{GF}(2^8)$ dan memakai
\textit{irreducible polynomial}
\[
f(x) = x^8 + x^4 + x^3 + x + 1.
\]

Sebagai contoh, misalkan $a = \code{57} = x^6 + x^4 + x^2 + x + 1$ dan
$b = \code{83} = x^7 + x + 1$. Hasil kali penuhnya adalah
\[
\begin{aligned}
a \cdot b
&= (x^6 + x^4 + x^2 + x + 1)(x^7 + x + 1) \\
&= x^{13} + x^{11} + x^9 + x^8 + 2x^7 + x^6 + x^5 + x^4 + x^3
+ 2x^2 + 2x + 1 \\
&\equiv x^{13} + x^{11} + x^9 + x^8 + x^6 + x^5 + x^4 + x^3 + 1
\pmod 2.
\end{aligned}
\]
Karena $m = 8$, hasil itu direduksi oleh $f(x)$:
\[
x^{13} + x^{11} + x^9 + x^8 + x^6 + x^5 + x^4 + x^3 + 1 \bmod f(x)
= x^7 + x^6 + 1 = \code{11000001}.
\]
Jadi $a \cdot b = \code{11000001} = \code{C1}$.

\subsection{Konstruksi S-box}

Elemen-elemen di dalam S-box tampak seperti diisi secara acak. Sebenarnya
elemen-elemen itu dihasilkan melalui perhitungan berikut.
\begin{enumerate}
    \item Inisialisasi S-box dengan nilai yang menaik dari baris ke baris. Baris
    ke-0 diisi dengan nilai $\code{00}, \code{01}, \code{02}, \dots, \code{0F}$.
    Baris ke-1 diisi dengan nilai $\code{10}, \code{11}, \code{12}, \dots,
    \code{1F}$, dan seterusnya untuk baris-baris lain.
    \item Untuk setiap nilai pada baris $y$ kolom $x$, tentukan balikannya di
    $\text{GF}(2^8)$. Nilai $\code{00}$ dipetakan ke dirinya sendiri.
    \item Hasil langkah 2 dikonversi ke vektor kolom bit
    $(b_0, b_1, \dots, b_7)^{\top}$.
    \item Kalikan vektor kolom bit itu dengan matriks \textit{affine} berikut.
    \begin{equation}
    \begin{bmatrix} b'_0 \\ b'_1 \\ b'_2 \\ b'_3 \\ b'_4 \\ b'_5 \\ b'_6 \\ b'_7 \end{bmatrix}
    =
    \begin{bmatrix}
    1 & 0 & 0 & 0 & 1 & 1 & 1 & 1 \\
    1 & 1 & 0 & 0 & 0 & 1 & 1 & 1 \\
    1 & 1 & 1 & 0 & 0 & 0 & 1 & 1 \\
    1 & 1 & 1 & 1 & 0 & 0 & 0 & 1 \\
    1 & 1 & 1 & 1 & 1 & 0 & 0 & 0 \\
    0 & 1 & 1 & 1 & 1 & 1 & 0 & 0 \\
    0 & 0 & 1 & 1 & 1 & 1 & 1 & 0 \\
    0 & 0 & 0 & 1 & 1 & 1 & 1 & 1
    \end{bmatrix}
    \begin{bmatrix} b_0 \\ b_1 \\ b_2 \\ b_3 \\ b_4 \\ b_5 \\ b_6 \\ b_7 \end{bmatrix}
    \oplus
    \begin{bmatrix} 1 \\ 1 \\ 0 \\ 0 \\ 0 \\ 1 \\ 1 \\ 0 \end{bmatrix}.
    \label{eq:aes-sbox-affine}
    \end{equation}
    \item Konversi hasil perhitungan $(b'_0, b'_1, \dots, b'_7)^{\top}$ ke dalam
    heksadesimal. Nilai itulah elemen $\text{S-box}(x,y)$.
\end{enumerate}

% Diagram lima langkah konstruksi S-box AES.
% Dirujuk dari section-04.tex dengan % Diagram lima langkah konstruksi S-box AES.
% Dirujuk dari section-04.tex dengan % Diagram lima langkah konstruksi S-box AES.
% Dirujuk dari section-04.tex dengan \input{figures/konstruksi-sbox}.
\begin{figure}[htbp]
    \centering
    \begin{tikzpicture}[
        kotak/.style={draw, rounded corners=1pt, minimum height=1.5cm,
                      text width=2.35cm, align=center, font=\small, inner sep=2pt},
        panah/.style={-{Latex[length=2mm]}, thick},
    ]
        \node[kotak] (a)  at (0,0)    {byte masukan\\$a = xy$};
        \node[kotak] (inv) at (2.9,0)  {invers di\\$\text{GF}(2^8)$\\$a^{-1}$};
        \node[kotak] (bit) at (5.8,0)  {vektor bit\\$(b_0,\dots,b_7)^\top$};
        \node[kotak] (aff) at (8.7,0)  {matriks affine\\$8\times8$\\$\oplus\ \code{63}$};
        \node[kotak] (out) at (11.6,0) {byte keluaran\\$\text{S-box}(x,y)$};

        \draw[panah] (a) -- (inv);
        \draw[panah] (inv) -- (bit);
        \draw[panah] (bit) -- (aff);
        \draw[panah] (aff) -- (out);
    \end{tikzpicture}
    \caption{Lima langkah konstruksi S-box AES. Byte masukan dibalik di $\text{GF}(2^8)$, hasilnya dinyatakan sebagai vektor bit, dikenai transformasi affine, lalu dibaca kembali sebagai byte}
    \label{fig:konstruksi-sbox}
    \sumber{Munir, \textit{IF4020 Kriptografi}, ITB.}
\end{figure}
.
\begin{figure}[htbp]
    \centering
    \begin{tikzpicture}[
        kotak/.style={draw, rounded corners=1pt, minimum height=1.5cm,
                      text width=2.35cm, align=center, font=\small, inner sep=2pt},
        panah/.style={-{Latex[length=2mm]}, thick},
    ]
        \node[kotak] (a)  at (0,0)    {byte masukan\\$a = xy$};
        \node[kotak] (inv) at (2.9,0)  {invers di\\$\text{GF}(2^8)$\\$a^{-1}$};
        \node[kotak] (bit) at (5.8,0)  {vektor bit\\$(b_0,\dots,b_7)^\top$};
        \node[kotak] (aff) at (8.7,0)  {matriks affine\\$8\times8$\\$\oplus\ \code{63}$};
        \node[kotak] (out) at (11.6,0) {byte keluaran\\$\text{S-box}(x,y)$};

        \draw[panah] (a) -- (inv);
        \draw[panah] (inv) -- (bit);
        \draw[panah] (bit) -- (aff);
        \draw[panah] (aff) -- (out);
    \end{tikzpicture}
    \caption{Lima langkah konstruksi S-box AES. Byte masukan dibalik di $\text{GF}(2^8)$, hasilnya dinyatakan sebagai vektor bit, dikenai transformasi affine, lalu dibaca kembali sebagai byte}
    \label{fig:konstruksi-sbox}
    \sumber{Munir, \textit{IF4020 Kriptografi}, ITB.}
\end{figure}
.
\begin{figure}[htbp]
    \centering
    \begin{tikzpicture}[
        kotak/.style={draw, rounded corners=1pt, minimum height=1.5cm,
                      text width=2.35cm, align=center, font=\small, inner sep=2pt},
        panah/.style={-{Latex[length=2mm]}, thick},
    ]
        \node[kotak] (a)  at (0,0)    {byte masukan\\$a = xy$};
        \node[kotak] (inv) at (2.9,0)  {invers di\\$\text{GF}(2^8)$\\$a^{-1}$};
        \node[kotak] (bit) at (5.8,0)  {vektor bit\\$(b_0,\dots,b_7)^\top$};
        \node[kotak] (aff) at (8.7,0)  {matriks affine\\$8\times8$\\$\oplus\ \code{63}$};
        \node[kotak] (out) at (11.6,0) {byte keluaran\\$\text{S-box}(x,y)$};

        \draw[panah] (a) -- (inv);
        \draw[panah] (inv) -- (bit);
        \draw[panah] (bit) -- (aff);
        \draw[panah] (aff) -- (out);
    \end{tikzpicture}
    \caption{Lima langkah konstruksi S-box AES. Byte masukan dibalik di $\text{GF}(2^8)$, hasilnya dinyatakan sebagai vektor bit, dikenai transformasi affine, lalu dibaca kembali sebagai byte}
    \label{fig:konstruksi-sbox}
    \sumber{Munir, \textit{IF4020 Kriptografi}, ITB.}
\end{figure}

